\chapter{Research Methodology}
\label{ch:methodology}

This appendix documents the research methodology used to produce this report, including the multi-agent research architecture, source selection criteria, codebase audit procedures, and Wardley mapping approach.

\section{Multi-Agent Research Architecture}

The research was conducted using five parallel research agents, each assigned a distinct domain with defined source priorities and quality thresholds. The agents operated concurrently across a shared memory layer (RuVector, 1.17 million vector embeddings with HNSW indexing) that enabled cross-domain pattern detection and prevented duplication of effort.

\begin{table}[H]
\centering
\caption{Research agent assignments and source priorities}
\label{tab:research-agents}
\begin{tabularx}{\textwidth}{l l X}
\toprule
\textbf{Agent} & \textbf{Domain} & \textbf{Primary Sources} \\
\midrule
Agent 1 & Mollick / BCG / Jagged Frontier & \textcite{DellAcqua2026}, \textcite{Randazzo2024}, \textcite{DellAcqua2023}, Mollick's ``One Useful Thing'' corpus, \textcite{BCGAIKPIs2025} \\
\addlinespace
Agent 2 & Block Critique & \textcite{DorseyBotha2026}, \textcite{Zamost2026}, \textcite{Spicer2026}, \textcite{Klarna2025}, \textcite{OrgvueForrester2025} \\
\addlinespace
Agent 3 & Agentic Organisations & \textcite{McKinseyAgenticOrg2026}, \textcite{McKinseySixShifts2026}, \textcite{BCGAgenticShift2025}, \textcite{BCGAgenticEnterprise2025}, \textcite{Deloitte2026} \\
\addlinespace
Agent 4 & Change Management & \textcite{WadeBonnet2025}, \textcite{Teece2025}, \textcite{MicrosoftDirectedAutonomy2025}, \textcite{Prosci2024}, \textcite{Capgemini2025} \\
\addlinespace
Agent 5 & Every / Compound Engineering & \textcite{Shipper2026}, \textcite{Klaassen2025}, \textcite{CyberArk2025}, \textcite{TDS2026} \\
\bottomrule
\end{tabularx}
\end{table}

Each agent maintained a source provenance log in the shared RuVector memory, recording the retrieval date, URL, and a confidence assessment for each source. Cross-referencing between agents was automated through semantic vector search: when Agent 2 (Block Critique) encountered a claim about AI productivity, the memory layer surfaced related findings from Agent 1 (Mollick / BCG) and Agent 5 (Every / Compound) for triangulation.

\section{Source Selection and Quality Assessment}

Sources were categorised into four tiers based on methodological rigour and evidentiary weight.

\subsection{Tier 1: Peer-Reviewed and Pre-Registered Research}

The highest evidentiary weight was assigned to peer-reviewed publications and pre-registered experiments. The primary load-bearing sources in this category are:

\begin{itemize}
  \item \textcite{DellAcqua2026}: pre-registered field experiment with 758 BCG consultants, published in \emph{Organization Science}. This is the evidential foundation for the jagged frontier concept.
  \item \textcite{Randazzo2024}: follow-up study with 244 BCG consultants identifying the Centaur, Cyborg, and Self-Automator collaboration modes.
  \item \textcite{DellAcqua2023}: field experiment with HR recruiters demonstrating the vigilance degradation problem.
  \item \textcite{Acemoglu2019}: the task-based automation framework that grounds the economic analysis of AI and labour.
  \item \textcite{Acemoglu2025}: the sceptical macroeconomic analysis of AI productivity, arguing for at most 0.66 per cent TFP gain over a decade.
\end{itemize}

\subsection{Tier 2: Institutional Research with Disclosed Methodology}

Consultancy reports and institutional research with disclosed sample sizes and methodologies. These carry less evidentiary weight than peer-reviewed work but are treated as credible within stated limitations:

\begin{itemize}
  \item \textcite{BCGAIKPIs2025}: joint research with MIT Sloan, $n = 3{,}000+$ across 25 industries.
  \item \textcite{MicrosoftWorkTrend2024}: large-scale survey with disclosed methodology.
  \item \textcite{WEF2025}: the World Economic Forum's Future of Jobs Report.
  \item \textcite{Capgemini2025}: middle manager resistance patterns with percentage-level findings.
  \item \textcite{EdmondsonLei2025}: psychological safety and AI experimentation, grounded in the established literature on team effectiveness.
\end{itemize}

\textbf{Limitation:} Consultancy reports are produced by organisations with commercial interests in the phenomena they describe. Methodology sections are often absent or sparse. Sample composition and response biases are rarely disclosed. Statistics from these sources are used directionally, to establish trends and orders of magnitude; rather than as precise measurements.

\subsection{Tier 3: Expert Commentary and Primary Accounts}

Expert commentary, practitioner accounts, and journalistic reporting. These provide context and narrative but are not treated as primary evidence:

\begin{itemize}
  \item \textcite{DorseyBotha2026}: primary source for the Block Intelligence Layer model.
  \item \textcite{Shipper2026} and \textcite{Klaassen2025}: primary accounts of the Every agent colony experiment.
  \item \textcite{Mollick2025Management}: expert synthesis from a researcher with deep domain expertise.
  \item \textcite{Spicer2026} and \textcite{Zamost2026}: expert and insider critique of Block.
\end{itemize}

\subsection{Tier 4: Community and Vendor Sources}

Sources with limited editorial oversight or commercial incentives. Used cautiously and flagged where they appear:

\begin{itemize}
  \item \textcite{TDS2026}: the original 17$\times$ error multiplication statistic. Superseded in this edition by \textcite{Cemri2025} (Chapter~\ref{ch:open-questions}); retained here in the bibliography for the correction record, not as an active source.
  \item \textcite{IBM2025} and \textcite{Gartner2025}: vendor-funded research with commercial incentives. Treated as directionally correct, per the methodological note in Chapter~\ref{ch:open-questions}.
\end{itemize}

\section{Web Search and Academic Database Approach}

Research was conducted through a combination of five source channels:

\begin{enumerate}
  \item \textbf{Academic databases}: Google Scholar, SSRN, NBER Working Papers, HBS Working Papers. Search terms included ``agentic organisation,'' ``AI coordination,'' ``human-AI collaboration,'' ``middle management AI,'' ``organisational design automation,'' and ``jagged technological frontier.''
  \item \textbf{Consultancy publication portals}: McKinsey Insights, BCG Publications, Deloitte Insights, Gartner Research. Filtered to 2024--2026 publications on organisational AI.
  \item \textbf{Practitioner platforms}: Substack (particularly Mollick's ``One Useful Thing'' and Every's ``Chain of Thought''), podcasts, and conference proceedings.
  \item \textbf{News and analysis}: Financial Times, New York Times, Fortune, CNBC, VentureBeat. Used for event-driven reporting (Block restructuring, Klarna reversal) with cross-referencing against multiple outlets.
  \item \textbf{Regulatory sources}: EUR-Lex for EU AI Act text, UK Legislation website for the Employment Rights Act, ICO guidance on workplace monitoring.
\end{enumerate}

\section{Codebase Audit Methodology}

The VisionClaw codebase was audited in April 2026 to reconcile claims made in earlier versions of this document against the actual code. The audit procedure:

\begin{enumerate}
  \item \textbf{Claim extraction}: all factual claims about VisionClaw's capabilities were extracted from the document and catalogued.
  \item \textbf{Code search}: each claim was tested against the codebase using structural code analysis (call graphs, dependency trees) and text search. The codebase-memory MCP tool (48,159 nodes, 95,766 edges) was used for structural queries.
  \item \textbf{Verification categories}: each claim was classified as (a) verified and accurate, (b) partially accurate with caveats, or (c) not supported by the code.
  \item \textbf{Correction}: claims in categories (b) and (c) were corrected in Chapter~\ref{ch:technical-substrate}. The corrections are presented transparently, with both the original claim and the corrected version visible.
\end{enumerate}

\begin{cautionbox}[title={Conflict of Interest}]
The codebase audit was conducted by the system architect (Dr.~John O'Hare), who is also the report author and the founder of DreamLab AI Consulting. This introduces a conflict of interest: the auditor has a commercial interest in the system being described favourably. The audit found and corrected seven overclaims (documented in Chapter~\ref{ch:technical-substrate}), which provides some evidence of good faith. Independent code review by a third party would strengthen the credibility of the technical claims but was not available at the time of writing.
\end{cautionbox}

\section{Wardley Mapping Methodology}

Five Wardley maps were produced for this report, each following Simon Wardley's standard mapping methodology:

\begin{enumerate}
  \item \textbf{Value chain identification}: for each map, the user need was identified and the value chain of components required to serve that need was enumerated.
  \item \textbf{Evolution assessment}: each component was placed on the evolution axis (Genesis $\rightarrow$ Custom $\rightarrow$ Product $\rightarrow$ Commodity/Utility) based on its current maturity and market availability.
  \item \textbf{Movement indication}: where components are expected to evolve (e.g., information routing moving from custom toward commodity), movement arrows indicate the direction and anticipated pace.
  \item \textbf{Strategic interpretation}: each map was interpreted for strategic implications relevant to the report's argument.
\end{enumerate}

The maps are interpretive tools, not empirical measurements. The placement of components on the evolution axis reflects the author's assessment informed by the evidence assembled in this report. Different analysts may position components differently, particularly for novel concepts (judgment brokerage, mesh velocity) whose evolutionary stage is inherently uncertain.

\section{The July 2026 Revision}

This edition was produced through a second, distinct research and drafting cycle three months after the first. The process:

\begin{enumerate}
  \item \textbf{Verification.} Four research agents, each with live web search, verified every transcript-derived claim in the April edition against primary sources --- the CAIS blog, arXiv, Bloomberg, the \emph{New York Times}, Ramp's own research publication, and court-reporting coverage of the Hangzhou ruling, correcting names, dates, and figures before any revision text was written. Corrections of substance are flagged explicitly at the point the original claim appears in the main text.
  \item \textbf{Ecosystem exploration.} Three further agents mapped the current state of the VisionClaw, AgentBox, and federation codebases to ground Chapter~\ref{ch:ecosystem-roadmap}, which applies this report's own findings to DreamLab's operational systems.
  \item \textbf{Drafting.} Chapter revisions were produced by section-scoped writing agents working against a single shared evidence brief with fixed citation keys, house style, and verified figures, so that no two chapters could independently introduce inconsistent numbers for the same claim. The drafts were then editorially reviewed for voice, argument continuity, and consistency with chapters outside each agent's assigned scope.
\end{enumerate}

The new sources introduced in this revision were tier-assigned using the same rubric as above:

\begin{itemize}
  \item \textbf{Tier 1}: \parencite{Cemri2025} --- peer-reviewed, NeurIPS 2025, Datasets and Benchmarks Track.
  \item \textbf{Tier 2}: \parencite{Liu2026} (preprint with disclosed methodology, single author); \parencite{Pereira2026Playbook} (disclosed methodology, acknowledged selection bias toward successful deployments); \parencite{RLI2025}, \parencite{CAIS2026RLI}, and \parencite{METR2026} (benchmark reports with disclosed methodology); \parencite{RampRevelio2026} (corporate research, disclosed methodology, correlational with matched controls).
  \item \textbf{Tier 3}: news reporting --- \parencite{Challenger2026}, \parencite{NYTChina2026}, \parencite{Bloomberg2026Ford}, \parencite{Chatterji2026}.
  \item \textbf{Tier 4}: \parencite{Levie2026}, a single social-media source with no published methodology, flagged wherever it is used.
\end{itemize}

\section{The Third Edition: the Gap-Analysis Mesh}

The third edition's principal new evidence artefact is the four-surface gap register (Chapter~\ref{ch:gap-register}). It was produced by a twelve-agent analysis mesh arranged in two layers, run on 8 July 2026, followed by a separate ten-agent drafting pass.

\textbf{Evidence layer (eight agents).} One agent extracted the ecosystem's espoused doctrine from the VisionFlow canon --- the claims the documentation actually commits to, with file-level citations. Four agents surveyed practice, one per interaction surface (forum, desktop, mixed reality, voice), verifying every capability claim against the running codebases and classifying each as working, partial, scaffolded, docs-only, or absent. Three agents researched external theory and the 2025--26 practice landscape through deliberately independent search channels: a closed synthesis engine, an open multi-source engine with citation verification against academic databases, and a high-recall keyword engine. The three theory agents did not coordinate and their briefs were identical; convergence across them was treated as a signal of consensus in the literature, divergence as a flag for editorial caution.

\textbf{Judgment layer (four agents).} One judge per surface crossed the three inputs --- external theory, espoused canon, verified practice, and recorded gaps typed by which pair diverged: theory-to-canon (published knowledge the doctrine never adopted), canon-to-practice (doctrine promises the code does not deliver), and theory-to-practice (absent from both). Judges were required to re-verify, by reading code, the claims their highest-severity findings depended on, and to check each candidate finding against the prior May 2026 audit so that already-known gaps were marked as still open rather than presented as new.

\textbf{Drafting.} The new chapters of Parts~\ref{part:solution-space}--\ref{part:next-steps} were then written by ten section-scoped writing agents working against a single binding style brief --- house voice, a prose-quality rule set, fixed citation keys, and an instruction hierarchy that ranks the code-verified surveys above aspirational documentation wherever they conflict. New bibliography entries were restricted to sources surfaced with checkable URLs during the theory research; a writer finding no source for a claim was instructed to state it as this report's own analysis or cut it. The drafts received a single-editor review for argument continuity and cross-chapter consistency.

\textbf{Ontology grounding.} The pervasive ontology binding described in Chapter~\ref{ch:ontology-binding} was active throughout production. Its fail-open contract was exercised in the field during this edition's production: the knowledge-graph backend was unreachable during one working session, every retrieval degraded to empty, and the work continued from documentation --- an observed instance of the design behaving as specified, recorded in that chapter.

\begin{cautionbox}[title={Limits of the Mesh Method}]
The analysis mesh shares a single foundation-model family. \textcite{Liu2026}'s pseudo-diversity result --- that same-model, same-tooling reviewers approximate one judge asked several times. This applies to this analysis as much as to any system it audits. The three search channels diversify retrieval, and the code-verification requirement anchors findings to artefacts rather than opinion, but the judgment layer is not model-diverse, and the auditor remains the audited organisation. The register's severity ratings should be read with both caveats attached.
\end{cautionbox}

\subsection{Weaving in the 2022 Baseline}

A separate tiered swarm produced the material that connects this report to the author's 2022 monograph \parencite{OHare2022Convergence}, disclosed in the introduction (Chapter~\ref{ch:introduction}) and audited in Chapter~\ref{ch:2022-wager}. It ran in four ranked stages, each narrower and more senior than the one before.

\textbf{Scouts (sixteen agents).} Sixteen Haiku agents read the 42 LaTeX source files of the 2022 book, a handful each, and extracted candidate continuity threads (passages on identity, Bitcoin, Lightning, decentralised organisations, mixed reality and language models that prefigured a 2026 component) with file-level citations and no ruling on strength.

\textbf{Synthesists (four agents).} Four Sonnet agents consolidated the scout output into a single lineage map, typing each thread by evidential strength: verbatim-level continuity, strong on mechanism but loose on narrative, and thematic only. Threads that could not be tied to a specific 2022 passage were dropped.

\textbf{Integrators (five agents).} Five Opus agents wove the surviving threads into the existing chapters as additive sentences, under the report's disclosure discipline: state the code-verified 2026 fact first, then name its 2022 ancestry, and retrofit no continuity that reads as marketing. Each integrator was scoped to a fixed set of files so that no two edited the same passage in parallel.

\textbf{Editor (single pass).} One editor reconciled voice across the insertions, cut overclaims, and required every continuity sentence to point forward to the chapter carrying the code-verified detail rather than restating it.

Two standing resources grounded the pass beyond the 2022 text. The operator's RuVector memory, through its \texttt{personal-context} and \texttt{project-state} namespaces, supplied the verified authorial history (the telepresence doctorate and the Octave mixed-reality laboratory tenure) that keeps the continuity claims accurate without inflating them. The pervasive ontology binding of Chapter~\ref{ch:ontology-binding}, active here as it was for the gap-analysis mesh, bound recurring terms such as the identity keypair, the block-trail and the judgment broker to their governed definitions, so the 2022-to-2026 mapping ran on one vocabulary at both ends.

\section{Limitations of the Research Process}

\begin{enumerate}
  \item \textbf{Single-author synthesis}: while the research was conducted using multiple agents, the synthesis and editorial judgment are those of a single author. The report reflects one analytical perspective on a multi-faceted phenomenon.
  \item \textbf{Temporal constraints}: this edition incorporates sources through early July 2026; sources published after that point are not included. The jagged frontier, in particular, continues to shift --- the Remote Labor Index update alone recorded a 6.4-fold increase in client-quality automation in the eight months preceding this revision \parencite{CAIS2026RLI}, and there is no reason to expect the pace to slow before the next edition.
  \item \textbf{English-language bias}: the research draws almost exclusively on English-language sources. Organisational innovations in non-Anglophone contexts, particularly Haier's micro-enterprise model in China, are discussed secondhand rather than from primary sources.
  \item \textbf{Analyst-vendor conflict}: the author is the founder of DreamLab AI Consulting and the architect of VisionClaw. This conflict is acknowledged throughout the report and addressed directly in Chapter~\ref{ch:open-questions}.
  \item \textbf{Sample bias in practitioner evidence}: the Every experiment ($\sim$20 people) and VisionClaw deployment ($\sim$15--50 people) are both small-scale, technically sophisticated organisations. The generalisability of findings from these contexts to larger, less technical organisations is an open question.
\end{enumerate}

\section{Disclosure}

The author of this report, Dr.\ John O'Hare, is the founder and principal of DreamLab AI Consulting Ltd and the architect of the VisionClaw system described in Chapter~\ref{ch:technical-substrate}. The report simultaneously analyses the problem space, proposes a solution framework, and describes a commercial offering from the author's firm. This positioning conflict is acknowledged in Chapter~\ref{ch:open-questions} and is offered as context for the reader's own evaluation of the argument.
